\begin{frame}{공식을 극한으로 — 그리고 $\rho$를 손으로}
  \begin{columns}[T]
    \begin{column}{0.5\textwidth}
      \textbf{두 극한}
      \begin{itemize}
        \item \textbf{완전히 독립}($\rho=0$): 분산 $\sigma^2/T$,
          표준편차 $\sigma/\sqrt{T}$로 \textbf{줄어든다}
        \item \textbf{완전히 같으면}($\rho=1$): 분산 $\sigma^2$
          \textbf{그대로} — ``완전히 같은 실수를 하는 전문가 100명''은
          전문가 1명과 \textbf{같다}
      \end{itemize}
      \vskip0.5em
      현실은 둘 사이($0<\rho<1$):
      \textbf{무작위성이 클수록} $\rho$가 작아져 $\sigma/\sqrt{T}$에
      가까워진다.
    \end{column}
    \begin{column}{0.46\textwidth}
      \textbf{손으로}
      \begin{itemize}
        \item 개별 트리 $\sigma \approx 1.55\%$p, 독립이라면 100개 평균은
          $\tfrac{1.55}{\sqrt{100}} = 0.155\%$p
        \item 실측: 5번 다른 시드 포레스트가 \textbf{0.68\%p}로 모였다
          — ``완전히 독립(0.155\%p)은 아니지만 개별 트리보다
          훨씬 안정''
        \item 0.68\%p를 공식에 대입하면
          $\rho \approx 0.2$ — 같은 강한 특징을 다 보지만 \textbf{상관이
          낮은 쪽}
      \end{itemize}
    \end{column}
  \end{columns}
  \vskip0.6em
  \begin{block}{결론}
    두 가지 무작위성의 역할 = 정확히 이 \textbf{$\rho$를 줄이는 것}.
    $\rho$가 1에 가깝지 않다는 것이 ``앙상블이 단일 트리보다 좋은 이유''의
    \textbf{전부}.
  \end{block}
\end{frame}
